% Part: computability
% Chapter: recursive-functions
% Section: composition

\documentclass[../../../include/open-logic-section]{subfiles}

\begin{document}

% SEGMENT OLP-0212-S01

\olfileid{cmp}{rec}{com}
\olsection{சார்புகளின் சேர்ப்பு}

$f$ மற்றும் $g$ இயல் எண்களின் இரு ஓரிடச் சார்புகள் எனில்,
அவற்றைச் சேர்க்கலாம்: $h(x) = f(g(x))$. அப்போது புதிய
சார்பு~$h(x)$, சார்புகள் $f$ மற்றும்~$g$ இலிருந்து
\emph{சேர்ப்பால்} வரையறுக்கப்படுகிறது. ஒன்றுக்கு மேற்பட்ட
உள்ளீடுகளைக் கொண்ட சார்புகளுக்கும் இதைப் பொதுமைப்படுத்த விரும்புகிறோம்.

இதைச் செய்வதற்கான ஒரு வழி பின்வருமாறு: $f$ ஒரு $k$-இடச் சார்பு
என்றும், $g_0$, \dots, $g_{k-1}$ ஆகியவை அனைத்தும் $n$-இடச்
சார்புகளான $k$ சார்புகள் என்றும் கொள்க. அப்போது புதிய $n$-இடச்
சார்பு $h$ ஐப் பின்வருமாறு வரையறுக்கலாம்:
\[
h(x_0, \dots, x_{n-1}) =
f(g_0(x_0, \dots, x_{n-1}), \dots, g_{k-1}(x_0, \dots, x_{n-1}))
\]
$f$ மற்றும் அனைத்து $g_i$ உம் கணிக்கத்தக்கவை எனில், $h$ உம்
கணிக்கத்தக்கது: $h(x_0, \dots, x_{n-1})$ ஐக் கணிக்க, முதலில்
ஒவ்வொரு $i = 0$, \dots,~$k-1$ க்கும் $y_i = g_i(x_0, \dots,
x_{n-1})$ என்ற மதிப்புகளைக் கணிக்க. பின்னர் இம்மதிப்புகளை $f$ க்கு
உள்ளீடாக அளித்து $h(x_0, \dots, x_{k-1}) = f(y_0, \dots, y_{k-1})$
ஐக் கணிக்க.

ஏற்கெனவே உள்ள சில சார்புகளைப் பயன்படுத்தி ஒரு புதிய சார்பைக்
கணிக்கும்போது நிகழ்வதை இது மிகக் கட்டுப்படுத்திய முறையில்
வகைப்படுத்துவதுபோல் தோன்றலாம். முதன்மை மீள்வரையறை வழி~$\Add$ ஐ
வரையறுப்பதில் $g(x, y, z) = \Succ(z)$ ஐப் பயன்படுத்தியதைப் போல,
சிலவேளைகளில் ஒரு சார்பின் எல்லா உள்ளீடுகளையும் நாம் பயன்படுத்துவதில்லை.
பின்வரும் சார்புகளைப் பயன்படுத்த நமக்கு அனுமதி உண்டு என்று கொள்க:
\[
\Proj{n}{i}(x_0, \dots, x_{n-1}) = x_i
\]
சார்புகள்~$\Proj{k}{i}$ \emph{வீழ்த்தல்} சார்புகள் எனப்படுகின்றன:
$\Proj{n}{i}$ ஓர் $n$-இடச் சார்பு. அப்போது $g$ ஐப் பின்வருமாறு
வரையறுக்கலாம்:
\[
g(x, y, z) = \Succ(\Proj{3}{2}(x, y, z)).
\]
இங்கே $f$ இன் பணியை $1$-இடச் சார்பு $\Succ$ வகிக்கிறது; எனவே
$k=1$. மேலும் $g_0$ இன் பணியை வகிக்கும் ஒரு $3$-இடச் சார்பு
$\Proj{3}{2}$ நம்மிடம் உள்ளது. விளைவு, மூன்றாம் உள்ளீட்டின்
அடுத்தெண்ணைத் தரும் $3$-இடச் சார்பு.

உள்ளீடுகளை மறுவரிசைப்படுத்தியோ ஒன்றாக அடையாளப்படுத்தியோ புதிய
சார்புகளை வரையறுக்கவும் வீழ்த்தல் சார்புகள் நம்மை அனுமதிக்கின்றன.
எடுத்துக்காட்டாக, சார்பு $h(x) = \Add(x, x)$ ஐப் பின்வருமாறு
வரையறுக்கலாம்:
\[
h(x_0) = \Add(\Proj{1}{0}(x_0),\Proj{1}{0}(x_0)).
\]
இங்கே $k=2$, $n=1$; $f(y_0,y_1)$ இன் பணியை $\Add$ வகிக்கிறது;
$g_0(x_0)$ மற்றும் $g_1(x_0)$ ஆகிய இரண்டின் பணிகளையும்
~$\Proj{1}{0}(x_0)$ என்ற ஓரிட வீழ்த்தல் சார்பு (அதாவது ஒன்றாதல்
சார்பு) வகிக்கிறது.

$f(y_0, y_1)$ என்பது ஏற்கெனவே நம்மிடம் உள்ள சார்பு எனில்,
சார்பு $h(x_0, x_1) = f(x_1, x_0)$ ஐப் பின்வருமாறு வரையறுக்கலாம்:
\[
h(x_0, x_1) = f(\Proj{2}{1}(x_0, x_1),\Proj{2}{0}(x_0, x_1)).
\]
இங்கே $k=2$, $n = 2$; $g_0$ மற்றும் $g_1$ இன் பணிகளை முறையே
$\Proj{2}{1}$ மற்றும்~$\Proj{2}{0}$ வகிக்கின்றன.

$g_0$, \dots,~$g_{k-1}$ அனைத்தும் ஒரே இடஎண்~$n$ ஐக்
கொண்டிருக்க வேண்டும் என்பதும் கவலையளிக்கலாம். (ஒரு சார்பின்
\emph{இடஎண்} என்பது அதன் உள்ளீடுகளின் எண்ணிக்கை என்பதை நினைவுகூர்க;
ஓர் $n$-இடச் சார்பின் இடஎண்~$n$.) ஆனால் வீழ்த்தல் சார்புகளைச்
சேர்ப்பது விரும்பிய நெகிழ்வைத் தருகிறது. எடுத்துக்காட்டாக,
$f$ மற்றும்~$g$ ஆகியவை $3$-இடச் சார்புகள் என்றும், பின்வருமாறு
வரையறுக்கப்படும் $2$-இடச் சார்பு~$h$ என்றும் கொள்க:
\[
h(x,y) = f(x,g(x,x,y),y).
\]
வீழ்த்தல் சார்புகளைக் கொண்டு~$h$ இன் வரையறையைப் பின்வருமாறு
மீண்டும் எழுதலாம்:
\[
h(x,y) = f(\Proj{2}{0}(x,y), g(\Proj{2}{0}(x,y), \Proj{2}{0}(x,y),
\Proj{2}{1}(x,y)), \Proj{2}{1}(x,y)).
\]
அப்போது $h$ ஆனது $f$ உடன் $\Proj{2}{0}$, $l$ மற்றும்
$\Proj{2}{1}$ ஆகியவற்றின் சேர்ப்பு; இங்கே
\[
l(x,y) = g(\Proj{2}{0}(x,y),\Proj{2}{0}(x,y),\Proj{2}{1}(x,y)),
\]
அதாவது $l$ ஆனது $g$ உடன் $\Proj{2}{0}$, $\Proj{2}{0}$ மற்றும்
~$\Proj{2}{1}$ ஆகியவற்றின் சேர்ப்பு.

\end{document}
